% !TeX root = ../初等数学笔记.tex
% 由主文档通过 \input 载入；请编译主文档。
\setcounter{chapter}{8} % 保留原书章号。
\chapter{数列、递推与求和}
\label{chap:em-sequences}
数列把函数的输入限制为整数指标，适合描述逐次累积、递推计算与离散变化。求通项和求和的共同思路，是寻找相邻项之间能够消去或简化的关系。

\section{数列、通项与递推关系}
\emph{数列}（sequence）是定义在正整数或其初始有限段上的函数，第 $n$ 项记为 $a_n$。前 $n$ 项和记为 $S_n=\sum_{k=1}^n a_k$，于是 $a_1=S_1$，而 $n\ge2$ 时 $a_n=S_n-S_{n-1}$。后一个关系不能在未定义 $S_0$ 时直接用于 $n=1$；若约定空和 $S_0=0$，则可统一写出。

递推关系用先前项决定后续项，必须配有足够初值。例如 $a_{n+1}=2a_n$ 配上不同 $a_1$ 就产生不同数列。有限个已知项也不能唯一决定无限数列：在某个候选通项 $f(n)$ 上增加 $c(n-1)\cdots(n-m)$，前 $m$ 项不变，后续通常改变。因此发现规律之后，还要从题设条件证明它。

\section{等差数列与等比数列}
若 $a_{n+1}-a_n=d$，则累加相邻差得 $a_n=a_1+(n-1)d$。将正序与倒序的和相加，每对均为 $a_1+a_n$，因此
\[
 S_n=\frac{n(a_1+a_n)}2=na_1+\frac{n(n-1)}2d.
\]
若 $a_{n+1}=qa_n$，且 $q\ne0$，则 $a_n=a_1q^{n-1}$。对任意实数 $q\ne1$，错位相减得到
\begin{equation}\label{eq:em-geometric-sum}
 (1-q)S_n=a_1(1-q^n),\qquad
 S_n=\frac{a_1(1-q^n)}{1-q}.
\end{equation}
$q=1$ 时 $S_n=na_1$；$q=0$ 时从第二项起全为零，$S_n=a_1$，直接由递推解释即可。零项情形不宜用 $a_{n+1}/a_n$ 定义公比。

\section{裂项、错位相减与分组求和}
\emph{裂项求和}（telescoping sum）利用 $a_k=b_k-b_{k+1}$ 使中间项抵消，例如
\[
 \sum_{k=1}^n\frac1{k(k+1)}
 =\sum_{k=1}^n\left(\frac1k-\frac1{k+1}\right)
 =1-\frac1{n+1}.
\]
间隔不是 $1$ 时须保留相应数量的首尾项，例如
\[
 \sum_{k=1}^n\frac1{k(k+2)}
 =\frac12\left(1+\frac12-\frac1{n+1}-\frac1{n+2}\right).
\]
根式也可裂项：$1/(\sqrt{k+1}+\sqrt k)=\sqrt{k+1}-\sqrt k$，前 $n$ 项和为 $\sqrt{n+1}-1$。

若 $T_n=\sum_{k=1}^n kq^{k-1}$，$q\ne0,1$，则 $T_n-qT_n=1+q+\cdots+q^{n-1}-nq^n$，结合式 \eqref{eq:em-geometric-sum} 得
\[
 T_n=\frac{1-(n+1)q^n+nq^{n+1}}{(1-q)^2}.
\]
这说明错位相减不仅能求等比和，还能降低等差因子的次数；$q=1$ 改用等差和，$q=0$ 则直接按原来的多项式求和约定计算。

\section{简单递推的化归}
对 $a_{n+1}=pa_n+q$，若 $p=1$，则为公差 $q$ 的等差数列；若 $p\ne1$，令 $L=q/(1-p)$，有 $a_{n+1}-L=p(a_n-L)$。当 $p\ne0,1$ 时
\[
 a_n=L+p^{n-1}(a_1-L).
\]
$p=0$ 时从第二项起恒为 $q$，不必引入 $0^0$。这里 $L$ 是更新规则的不动点，平移把非齐次递推化为等比递推。

二阶例子 $F_1=F_2=1,F_{n+2}=F_{n+1}+F_n$ 的候选指数解满足 $r^2=r+1$。记 $\alpha=(1+\sqrt5)/2,\beta=(1-\sqrt5)/2$，直接代入初值与递推可验证 $F_n=(\alpha^n-\beta^n)/\sqrt5$。唯一性由递推逐项确定保证，因此验证完整初值和递推已经足以证明公式。

\section{单调、有界与极限的预备认识}
若对每个 $n$ 有 $a_{n+1}\ge a_n$，数列非严格递增；若还有统一上界，就能用完备性证明它收敛。这里 $a_n\to L$ 的严格含义是：对每个 $\varepsilon>0$，存在正整数 $N$，使所有 $n\ge N$ 都满足 $|a_n-L|<\varepsilon$。例如 $1/n\to0$，因为取整数 $N>1/\varepsilon$ 即可。

\begin{theorem}[单调有界收敛]
递增且有上界的实数列收敛到其所有项的上确界；递减且有下界的情形相应成立。
\end{theorem}
\begin{proof}
设 $L=\sup\{a_n:n\ge1\}$。对 $\varepsilon>0$，$L-\varepsilon$ 不是上界，故存在 $N$ 使 $a_N>L-\varepsilon$。递增性给出 $L-\varepsilon<a_N\le a_n\le L$，对所有 $n\ge N$ 成立，这就是极限定义。递减情形对 $-a_n$ 应用已证结果。
\end{proof}
当 $|q|<1$ 时 $q^n\to0$：$q=0$ 直接成立；否则令 $u=1/|q|-1>0$，归纳得 $(1+u)^n\ge1+nu$，所以 $|q|^n\le1/(1+nu)\to0$。因此等比级数的部分和在该条件下趋于 $a_1/(1-q)$；没有收敛条件，有限和公式不能直接变成无穷和。

\section{例题、习题与学习检查}
\begin{example}
设 $a_1=1,a_{n+1}=2a_n+3$，求通项。
\end{example}
\begin{solve}
不动点为 $L=-3$，故 $a_{n+1}+3=2(a_n+3)$，于是 $a_n+3=4\cdot2^{n-1}$，即 $a_n=2^{n+1}-3$。代入 $n=1$ 及更新式均吻合，确认初值和指标没有错位。
\end{solve}
\begin{enumerate}
\item 求 $\sum_{k=1}^n(2k-1)$ 与 $\sum_{k=1}^n3^{k-1}$。答：$n^2$ 与 $(3^n-1)/2$。
\item 设 $a_1=0,a_{n+1}=(a_n+2)/2$，求通项和极限。答：$a_n=2-2^{2-n}$，极限为 $2$。
\item 证明 $a_n=1-1/n$ 严格递增且极限为 $1$，但任何一项都不等于极限。
\end{enumerate}

